% Input-only TikZ fragment.
% Main-document preamble requirements:
% \usepackage{graphicx,float,tikz}
% \usetikzlibrary{positioning,arrows.meta,calc}

\definecolor{biasred}{RGB}{175,35,35}
\definecolor{varteal}{RGB}{25,145,150}
\definecolor{totalblack}{RGB}{20,20,20}
\definecolor{boxbg}{RGB}{252,252,252}
\definecolor{ph1}{RGB}{170,35,35}
\definecolor{ph2}{RGB}{45,110,55}
\definecolor{ph3}{RGB}{30,95,165}

\begin{tikzpicture}[font=\normalsize]

% ---------- geometry ----------
\def\xmax{7.2}
\def\ymax{5.3}
\def\xopt{4.13}
\def\yopt{0.85}
\def\timex{7.95}
\def\boxx{8.75}

% ---------- styles ----------
\tikzset{
  ax/.style={-{Latex[length=3mm,width=2.3mm]}, line width=1.3pt, totalblack},
  cv/.style={line width=2.1pt, samples=220, domain=0.02:9.5},
  lead/.style={line width=0.6pt, opacity=0.85},
  phasebox/.style={draw=#1, fill=boxbg, rounded corners=5pt, line width=1pt,
                   text width=5.15cm, inner xsep=6pt, inner ysep=5pt, % was 6pt
                   align=left, font=\small},
}

% Shift the complete landscape left; the timeline and phase boxes stay fixed.
\begin{scope}[xshift=-5mm]

% ---------- axes ----------
\draw[ax] (-0.2,0) -- (\xmax+0.25,0);
\draw[ax] (0,-0.2) -- (0,\ymax+0.25);
\node[font=\large\bfseries, totalblack, anchor=north west]
      at (5.10,-0.24) {Complexity};
\node[font=\large\bfseries, totalblack, rotate=90]
      at (-0.58,{0.46*\ymax}) {Error};

% ---------- curves ----------
\begin{scope}
  \clip (0,0) rectangle (\xmax,\ymax);
  \draw[cv,biasred]
    plot (\x,{5*exp(-0.6*\x)+0.07});
  \draw[cv,varteal]
    plot (\x,{0.02*exp(0.7*\x)});
  \draw[cv,totalblack,line width=2.4pt]
    plot (\x,{5*exp(-0.6*\x)+0.02*exp(0.7*\x)+0.07});
\end{scope}

% ---------- optimum ----------
\draw[totalblack!45, dashed, line width=1pt]
  (\xopt,0.08) -- (\xopt,\ymax-0.35);
\filldraw[fill=white, draw=totalblack, line width=1pt]
  (\xopt,0) circle (0.14);
\fill[totalblack] (\xopt,\yopt) circle (0.07);
\node[font=\small\bfseries, totalblack!75, anchor=north east]
  at (\xopt-0.12,-0.25) {Optimum};
\node[font=\footnotesize\itshape, totalblack!60, anchor=north,
      align=center, text width=4.8cm]
  at (\xopt,-0.77)
  {Position of current practice relative to this optimum unknown};

% ---------- curve labels ----------
\node[biasred, font=\normalsize\bfseries, anchor=west]
  at (1.00,0.80) {Bias$^2$};
\draw[lead,biasred] (2.18,0.93) -- (2.53,1.12);
\fill[biasred] (2.55,1.15) circle (0.05);

\node[varteal, font=\normalsize\bfseries, anchor=west]
  at (5.35,0.55) {Variance};
\draw[lead,varteal] (5.95,0.78) -- (6.17,1.46);
\fill[varteal] (6.20,1.53) circle (0.05);

\node[totalblack, font=\normalsize\bfseries, anchor=east]
  at (3.90,3.00) {Total Error};
\draw[lead,totalblack] (3.45,2.80) -- (3.07,1.11);
\fill[totalblack] (3.05,1.04) circle (0.05);

% ---------- unknown current position ----------
% The two free-floating "?" glyphs are gone: they read as stray characters.
% The "position undetermined" statement now lives in the note under the axis,
% in the Phase III box, and in the caption. If you want a visual hook back,
% use this self-explanatory badge instead of bare question marks:
% \node[draw=ph3!50, dashed, thick, rounded corners=4pt, align=center,
%       inner sep=4pt, text=ph3!75, font=\large\bfseries]
%   at (5.90,4.20) {?\\[-1pt]\scriptsize\normalfont\itshape current practice};

\end{scope}

% ---------- phase boxes (tighter column: 6 mm gaps, 5 pt y-sep) ----------
\node[phasebox=ph1, anchor=north west] (box1) at (\boxx,\ymax)
  {\textcolor{ph1}{\large\bfseries Phase I}\hspace{5pt}%
   \textcolor{ph1}{\footnotesize Pre-2005}\\[4pt]
   Siloed models\\[1pt]
   Disciplinary silos};

\node[phasebox=ph2, below=6mm of box1] (box2)
  {\textcolor{ph2}{\large\bfseries Phase II}\hspace{5pt}%
   \textcolor{ph2}{\footnotesize 2005--?}\\[4pt]
   Institutional integration\\[1pt]
   Participatory approaches};

\node[phasebox=ph3, below=6mm of box2] (box3)
  {\textcolor{ph3}{\large\bfseries Phase III}\hspace{5pt}%
   \textcolor{ph3}{\footnotesize ? (unverified)}\\[4pt]
   ``Model of everything''\\[1pt]
   Position on landscape \textbf{undetermined}};

\draw[-{Latex[length=2.2mm,width=1.8mm]}, black!35, line width=0.9pt]
  (box1) -- (box2);
\draw[-{Latex[length=2.2mm,width=1.8mm]}, black!35, line width=0.9pt]
  (box2) -- (box3);

% ---------- temporal separator (auto-centred on the box column) ----------
\coordinate (timelinetop) at (\timex,\ymax);
\coordinate (timelinebottom) at (timelinetop |- box3.south);
\draw[-{Latex[length=3mm,width=2.3mm]}, black!40, line width=1.2pt]
  (timelinetop) -- (timelinebottom);
\node[font=\small\itshape, black!55, rotate=90, anchor=center]
  at ([xshift=-4mm]$(timelinetop)!0.5!(timelinebottom)$)
  {Temporal evolution};

\end{tikzpicture}